% !TeX root = ../papers/nc-p3.tex
\section{核直線族の非被約因子と四次元中心の三次分岐の排除}
\label{sec:rank-ten-cubic-kernel-divisor}
\begin{theorem}\label{thm:rank-ten-cubic-kernel-divisor}
反対階数10で曲面射影のない標準正則窓の中央支持4の順中心の第二曲線は、
必ず滑らかな平面二次曲線である。第一直線上の唯一の順中心は特異三項鉛筆型である。
第三曲線および次の位置は定理\ref{thm:rank-ten-conic-after-center}の二平面に入る。
逆向きにも成立する。
\end{theorem}
\begin{proof}[核直線上での次辺の必要条件]
第二曲線がねじれ三次であると仮定する。定理\ref{thm:rank-ten-cubic-after-center}の
I, II型のいずれでも、極限を含む核直線を $t=(s,t')$ と $[X:Y]$ の積でパラメータ化する。
混同を避けて以下 $t=(t_0,t_1)$ とし、位置3の座標をそれぞれ
\[
 z_{\rm I}=(-t_1X,t_0X,-t_1Y,t_0Y),\qquad
 z_{\rm II}=(-t_1X,t_0X-t_1Y,t_0Y,0)
\]
とする。いずれも基点のない $(1,1)$ 式であり、$A(t)z=0$。
したがって $E(t)z=w(-,-,x_2(t),z)\in R_0$ は第一直線への制限写像の
固定一次元核に値を取る。その係数は一つの $(4,1)$ 式 $D$ である。
恒等零なら全 $x_2(t)$ が順中心になり不可能である。
$D=0$ は二点 $x_2(t),z$ が実際に辺をなす条件そのもので、各固定 $t$ で $X,Y$ に関して一次である。
従って因子は垂直なファイバー（重複を含む）と、第三点のパラメータを与える重複度1の切断からなる。

この因子の各既約成分の一般点で、その重複度まで含めて $z$ に後続点がある。
実際積曲面の当該素因子での離散付値環を $D$ で割る。水平成分なら一般の $x_2(t)$ は
逆支持3、垂直成分なら $x_2(t)$ は順中心なのでやはり逆支持3である。
よってこの局所商環上でも逆評価の三列が六空間に階数3の直和因子を張る。
$w(-,x_2(t),z,-)=0$（巡回順での二点の辺の式）により、隣接切断の完全対合で
$z$ の順評価像はその直交補の階数3の直和因子に入る。
剰余体上で $z$ は動く一般点であり順支持3なので、単元な三次小行列で解けば
この局所商環で核直線が存在する。これが求める後続 $x_4$ である。
\end{proof}
\begin{proof}[次辺の小行列式の二重次数]
$R_3$ の中に三つの独立な行列 $A(y)$ の係数があり、残る三関係を $B_j$ とする。
I型で $Y$ が単元なら、全 $A(y)z$ との対合が消えるための $x_4$ の条件は、
前節の式より $v_1(x_4)=0$, $Xu(x_4)+Yv_0(x_4)=0$ である。
従って $x_4$ は $p$ と $[u:v_0:v_1]=[Y:-X:0]$ の固定線型持上げ $h(X,Y)$ の包に入る。
この二列はどの $[X:Y]$ でも独立である。残る三関係のこの包への制限行列は
$B_j(z,p)$ と $B_j(z,h)$ を二列に持つので、二次小行列式の二重次数は $(2,3)$ である。
前段の後続の存在により $Y=0$ の外で各小行列式は $D$ で割り切れる
（各素因子の重複度まで消える）。
$D$ の中の $Y$ 因子を全て取り除いても $t$ 次数は4なので、$t$ 次数2の小行列式は
全て恒等零である。こうして全ての一般の $z_{\rm I}$ が後続を持つ。
この $z$ の像は滑らかな二次曲面全体であり、曲面射影の不在に反する。

II型では $A(y)z$ は $u$ の二次式倍で、その全係数の一つは各 $z$ で単元である。
従って $x_4$ は平面 $u=0$ に入る。$B_j$ のその平面への制限を三行三列にすると、
行列式は $(3,3)$ 式である。これも $D$ の全素因子で重複度まで消えるので恒等零となる。
$z_{\rm II}$ の像は第四座標が零の全平面である。実際任意の二項二次式が根 $t$ を持ち、
その $t$ で表示の直線上に乗る。この全平面が後続を持つことになり矛盾する。
以上で第二次数3を排除した。二直線列はなく、残る二次曲線には既出の判定を適用できる。
\end{proof}
